% ============================================================
%  MANUAL DE USUARIO — CHATBOT FACULTAD
%  Facultad de Inteligencia Artificial e Ingenierías
%  Universidad de Caldas
%  Compilar con: pdflatex manual-usuario.tex  (×2)
% ============================================================

\documentclass[11pt, a4paper, oneside]{report}

% ── Codificación y tipografía ───────────────────────────────
\usepackage[utf8]{inputenc}
\usepackage[T1]{fontenc}
\usepackage[spanish,es-tabla]{babel}
\usepackage{lmodern}
\usepackage{microtype}

% ── Márgenes ────────────────────────────────────────────────
\usepackage[
  top=2.4cm, bottom=2.4cm,
  left=2.6cm, right=2.4cm,
  headheight=16pt
]{geometry}

% ── Colores institucionales ─────────────────────────────────
\usepackage[table, dvipsnames]{xcolor}
\definecolor{navy}{HTML}{00407D}
\definecolor{navydark}{HTML}{002D59}
\definecolor{navysoft}{HTML}{EDF3F9}
\definecolor{naranja}{HTML}{F27022}
\definecolor{naranjasoft}{HTML}{FFF4EC}
\definecolor{grisclaro}{HTML}{F5F7FA}
\definecolor{gristexto}{HTML}{333333}
\definecolor{grismuted}{HTML}{667689}
\definecolor{okgreen}{HTML}{027A48}
\definecolor{oksoft}{HTML}{ECFDF3}

% ── Fuente ──────────────────────────────────────────────────
\usepackage{helvet}
\renewcommand{\familydefault}{\sfdefault}

% ── Gráficos, tablas, cajas ─────────────────────────────────
\usepackage{graphicx}
\usepackage{float}
\usepackage{booktabs}
\usepackage{tabularx}
\usepackage{longtable}
\usepackage{array}
\usepackage{multirow}
\usepackage{tikz}
\usetikzlibrary{calc, positioning, shapes.geometric}
\usepackage[most]{tcolorbox}
\usepackage{fancyhdr}
\usepackage{enumitem}
\usepackage{titlesec}
\usepackage{tocbibind}
\usepackage{xcolor}
\usepackage[
  colorlinks=true,
  linkcolor=navy,
  urlcolor=naranja,
  citecolor=navy,
  pdftitle={Manual de Usuario — Chatbot Facultad},
  pdfauthor={Facultad de Inteligencia Artificial e Ingenierías — Universidad de Caldas},
  pdfsubject={Manual de Usuario del Asistente Virtual Institucional}
]{hyperref}

\newcolumntype{L}[1]{>{\raggedright\arraybackslash}p{#1}}
\newcolumntype{C}[1]{>{\centering\arraybackslash}p{#1}}

\setcounter{tocdepth}{2}
\setlength{\parskip}{6pt}
\setlength{\parindent}{0pt}

% ── Cajas ───────────────────────────────────────────────────
\newtcolorbox{notabox}{
  colback=navysoft, colframe=navy,
  fonttitle=\bfseries\color{navy},
  title={Nota},
  arc=5pt, boxrule=0.8pt,
  left=8pt, right=8pt, top=5pt, bottom=5pt
}

\newtcolorbox{advertenciabox}{
  colback=naranjasoft, colframe=naranja,
  fonttitle=\bfseries\color{naranja},
  title={Advertencia},
  arc=5pt, boxrule=0.8pt,
  left=8pt, right=8pt, top=5pt, bottom=5pt
}

\newtcolorbox{tipobox}{
  colback=oksoft, colframe=okgreen,
  fonttitle=\bfseries\color{okgreen},
  title={Recomendación},
  arc=5pt, boxrule=0.8pt,
  left=8pt, right=8pt, top=5pt, bottom=5pt
}

\newtcolorbox{pasobox}[1]{
  colback=grisclaro, colframe=navy,
  fonttitle=\bfseries\color{white},
  title={Paso #1},
  coltitle=white,
  attach boxed title to top left={yshift=-2.2mm, xshift=5mm},
  boxed title style={colback=navy, arc=3pt},
  arc=5pt, left=8pt, right=8pt, top=12pt, bottom=6pt
}

% ── Encabezado / pie ────────────────────────────────────────
\pagestyle{fancy}
\fancyhf{}
\fancyhead[L]{\small\textcolor{navy}{\nouppercase{\leftmark}}}
\fancyhead[R]{\small\textcolor{naranja}{\textbf{Chatbot Facultad}}}
\fancyfoot[C]{\small\textcolor{grismuted}{\thepage}}
\renewcommand{\headrulewidth}{1.2pt}
\renewcommand{\headrule}{%
  \hbox to\headwidth{%
    \color{navy}\leaders\hrule height 1.2pt\hfill
    \color{naranja}\vrule width 2.2cm height 1.2pt
  }%
}

% ── Títulos ─────────────────────────────────────────────────
\titleformat{\chapter}[display]
  {\normalfont}
  {\textcolor{naranja}{\rule{1.8cm}{3.5pt}}\\[0.35cm]
   {\large\bfseries\textcolor{navy}{Capítulo \thechapter}}}
  {0.35cm}
  {\Huge\bfseries\textcolor{navy}}
  [\vspace{0.25cm}{\color{naranja}\titlerule[2pt]}\vspace{0.15cm}{\color{navy!40}\titlerule[0.4pt]}]
\titlespacing*{\chapter}{0pt}{-12pt}{18pt}

\titleformat{\section}
  {\color{navy}\Large\bfseries}
  {\thesection}{0.55em}{}
  [\vspace{-2pt}{\color{naranja}\titlerule[1.2pt]}\vspace{4pt}]

\titleformat{\subsection}
  {\color{navy}\large\bfseries}
  {\thesubsection}{0.45em}{}

\titleformat{\subsubsection}
  {\color{gristexto}\normalsize\bfseries}
  {\thesubsubsection}{0.4em}{}

% ── Helpers UI ──────────────────────────────────────────────
\newcommand{\menupath}[1]{\textcolor{navy}{\textbf{#1}}}
\newcommand{\boton}[1]{%
  \colorbox{naranja}{\small\color{white}\bfseries\ #1\ }%
}
\newcommand{\botonnavy}[1]{%
  \colorbox{navy}{\small\color{white}\bfseries\ #1\ }%
}
\newcommand{\campo}[1]{\textit{\textcolor{gristexto}{#1}}}
\newcommand{\badgeok}[1]{%
  \colorbox{oksoft}{\small\color{okgreen}\bfseries\ #1\ }%
}
\newcommand{\badgeoff}[1]{%
  \colorbox{grisclaro}{\small\color{grismuted}\bfseries\ #1\ }%
}

% Placeholder de captura (reemplazar por \includegraphics cuando exista el PNG)
\newcommand{\captura}[2]{%
  % #1 = nombre archivo sugerido (sin extensión), #2 = descripción
  \begin{figure}[H]
    \centering
    \begin{tcolorbox}[
      enhanced,
      colback=navysoft,
      colframe=navy,
      boxrule=1pt,
      arc=8pt,
      width=0.94\linewidth,
      height=6.2cm,
      valign=center,
      halign=center,
      shadow={1mm}{-1mm}{0mm}{black!12},
      borderline west={4pt}{0pt}{naranja}
    ]
      {\large\textcolor{navy}{\bfseries[ Espacio para captura de pantalla ]}}\\[0.55cm]
      {\normalsize\textcolor{gristexto}{#2}}\\[0.7cm]
      {\footnotesize\textcolor{grismuted}{Coloque la imagen en:}}\\[0.15cm]
      {\footnotesize\texttt{documentation/img/#1.png}}\\[0.45cm]
      {\scriptsize\textcolor{naranja}{%
        Luego reemplace este bloque por:\\
        \texttt{\textbackslash includegraphics[width=0.92\textbackslash linewidth]\{img/#1.png\}}
      }}
    \end{tcolorbox}
    \caption{#2}
    \label{fig:#1}
  \end{figure}%
}

% ============================================================
\begin{document}

% ───────────────────────── PORTADA ─────────────────────────
\begin{titlepage}
  \newgeometry{margin=0cm}
  \begin{tikzpicture}[remember picture, overlay]
    % Fondo navy
    \fill[navy] (current page.south west) rectangle (current page.north east);
    % Franja inferior en azul más claro
    \fill[navy!70!white] ([yshift=0cm]current page.south west)
      rectangle ([yshift=1.6cm]current page.south east);
    % Franja diagonal decorativa (azul claro)
    \fill[white, opacity=0.08]
      ([yshift=6cm]current page.south west) --
      ([yshift=9.5cm]current page.south west) --
      ([yshift=4cm]current page.south east) --
      ([yshift=0.5cm]current page.south east) -- cycle;
    % Barra lateral azul
    \fill[navydark] (current page.north west)
      rectangle ([xshift=0.55cm]current page.south west);
  \end{tikzpicture}
  
  \color{white}
  \centering
  
  % Logo sin fondo blanco (transparente)
  \vspace*{1.8cm}
  \includegraphics[height=2.2cm]{logo-facultad1.png}
  \vspace{1.0cm}
  
  {\Large \bfseries Facultad de Inteligencia Artificial e Ingenierías}\\[0.15cm]
  {\large Universidad de Caldas}\\[2.0cm]
  
  {\color{naranja}\rule{0.5\linewidth}{2pt}}\\[0.6cm]
  
  {\fontsize{36}{42}\selectfont\bfseries Asistente Virtual}\\[0.1cm]
  {\fontsize{36}{42}\selectfont\bfseries Institucional}\\[0.6cm]
  
  {\LARGE \bfseries Manual de Usuario}\\[1cm]
  
  {\large Daniela Agudelo Martínez}\\[1cm]
  {\normalsize Manizales, Colombia}
  
  \vfill
  
  {\small Versión 1.0 \quad\textcolor{naranja}{|}\quad \today}\\[2.0cm]
  
  \restoregeometry
\end{titlepage}

\color{gristexto}
\pagecolor{white}

% ───────────────── Control de versiones ────────────────────
\newpage
\thispagestyle{empty}
\section*{Control de versiones}
\vspace{0.3cm}
\begin{tabularx}{\linewidth}{C{1.6cm} C{2.8cm} L{3.8cm} X}
  \toprule
  \rowcolor{navy}
  \textcolor{white}{\textbf{Versión}} &
  \textcolor{white}{\textbf{Fecha}} &
  \textcolor{white}{\textbf{Autor}} &
  \textcolor{white}{\textbf{Descripción}} \\
  \midrule
  1.0 & \today & Daniela Agudelo Martínez & Versión inicial del manual de usuario \\
  \bottomrule
\end{tabularx}

\vspace{1.2cm}
\begin{notabox}
  Este documento está dirigido al personal administrativo y técnico de la Facultad
  que opera el \textbf{panel de administración} del chatbot
\end{notabox}

\tableofcontents
\newpage
\listoffigures

% ============================================================
\chapter{Introducción}

\section{Propósito del manual}
Este manual describe el uso del \textbf{Asistente Virtual Institucional}
desarrollado para la Facultad de Inteligencia
Artificial e Ingenierías de la Universidad de Caldas.

El sistema permite:
\begin{itemize}
  \item Consultar documentos normativos e institucionales mediante un chat en el sitio web.
  \item Administrar el corpus documental, el modelo de inteligencia artificial y la apariencia del panel.
  \item Probar respuestas antes de publicarlas al público.
\end{itemize}

\section{¿A quién está dirigido?}
\begin{tabularx}{\linewidth}{L{4.2cm} X}
  \toprule
  \rowcolor{navy}
  \textcolor{white}{\textbf{Perfil}} &
  \textcolor{white}{\textbf{Uso principal}} \\
  \midrule
  Administrador del panel & Subir documentos, reindexar, configurar IA, gestionar usuarios. \\
  \rowcolor{navysoft}
  Personal de soporte & Orientar el uso del widget y verificar respuestas. \\
  Estudiantes / visitantes & Usar el chat flotante en el portal web de la Facultad. \\
  \bottomrule
\end{tabularx}

\section{Componentes del sistema}
El producto tiene tres partes visibles:

\begin{enumerate}
  \item \textbf{Widget de chat} en el sitio WordPress de la Facultad (uso público).
  \item \textbf{Panel de administración} web (uso restringido a administradores).
  \item \textbf{API} del motor RAG (backend; no se usa directamente en el día a día).
\end{enumerate}

\section{Convenciones de este manual}
\begin{tabularx}{\linewidth}{C{3.2cm} X}
  \toprule
  \rowcolor{navy}
  \textcolor{white}{\textbf{Convención}} &
  \textcolor{white}{\textbf{Significado}} \\
  \midrule
  \boton{Botón} & Acción principal (naranja). \\
  \rowcolor{navysoft}
  \botonnavy{Botón} & Acción secundaria (azul institucional). \\
  \campo{Campo} & Campo de formulario. \\
  \rowcolor{navysoft}
  \menupath{Menú} & Opción del menú lateral del panel. \\
  \badgeok{Activo} & Estado positivo o configurado. \\
  \bottomrule
\end{tabularx}

% ============================================================
\chapter{Widget público (sitio web)}

\section{Descripción}
El widget es un chat flotante que aparece en las páginas del sitio WordPress
de la Facultad. Permite hacer preguntas sobre documentos institucionales
indexados (normativa, programas, etc.).

\begin{notabox}
  El asistente responde \textbf{solo} con información recuperada de los documentos
  cargados. Si no encuentra la respuesta, indica que no tiene esa información.
\end{notabox}

\section{Cómo usarlo (estudiante o visitante)}

\begin{pasobox}{1}
  Abra el sitio web de la Facultad. En la esquina inferior (generalmente derecha)
  verá el botón flotante del asistente.
\end{pasobox}

\begin{pasobox}{2}
  Haga clic en el botón para abrir el panel de chat. La primera vez se mostrará
  el aviso de privacidad: léalo y acéptelo para continuar.
\end{pasobox}

\begin{pasobox}{3}
  Escriba su pregunta en el cuadro de texto y envíela. Espere la respuesta del
  asistente. Puede continuar con preguntas de seguimiento en la misma sesión.
\end{pasobox}

\begin{figure}[H]
  \centering
  \includegraphics[width=0.92\linewidth]{04-1widget-conversacion.png}
  \caption{Conversación de ejemplo}
  \label{fig:documentos-subida}
\end{figure}

\begin{figure}[H]
  \centering
  \includegraphics[width=0.92\linewidth]{04-2widget-conversacion.png}
  \caption{Conversación de ejemplo con fuentes citadas}
  \label{fig:documentos-subida}
\end{figure}

\begin{tipobox}
  Formule preguntas concretas, por ejemplo:
  \textit{«¿Quién es la segunda instancia para cancelar una asignatura?»}
  en lugar de preguntas demasiado generales.
\end{tipobox}

\section{Qué puede esperar el usuario}
\begin{itemize}
  \item Respuestas basadas en documentos oficiales indexados.
  \item Listado de fuentes (archivo / página) cuando el sistema las proporciona.
  \item Mensaje amigable de error si el servicio no está disponible temporalmente.
\end{itemize}

\begin{advertenciabox}
  El widget no sustituye trámites oficiales ni asesoría jurídica. Ante dudas
  críticas, confirme siempre con la dependencia correspondiente de la Facultad.
\end{advertenciabox}

% ============================================================
\chapter{Acceso al panel de administración}

\section{Requisitos}
\begin{itemize}
  \item Navegador actualizado (Chrome, Firefox, Edge o Safari).
  \item URL del panel proporcionada por el equipo técnico.
  \item Correo y contraseña de administrador asignados por la Facultad.
\end{itemize}

\section{Inicio de sesión}

\begin{pasobox}{1}
  Abra la URL del panel. Verá la pantalla de acceso con el branding de la Facultad.
\end{pasobox}

\begin{pasobox}{2}
  Complete \campo{Correo electrónico} y \campo{Contraseña} (mínimo 8 caracteres).
\end{pasobox}
\begin{pasobox}{3}
  Pulse \boton{Entrar}. Si las credenciales son correctas, accederá al panel.
\end{pasobox}

\begin{figure}[H]
  \centering
  \includegraphics[width=0.92\linewidth]{login-administracion.png}
  \caption{Pantalla de inicio de sesión del panel}
  \label{fig:documentos-subida}
\end{figure}

\begin{advertenciabox}
  Si el usuario está \textbf{inactivo}, el acceso será denegado aunque la
  contraseña sea correcta. Contacte a otro administrador para reactivarlo.
\end{advertenciabox}

\section{Navegación general}
Tras el login, la interfaz se organiza así:

\begin{itemize}
  \item \textbf{Barra lateral (izquierda):} logo, menú (Inicio, Documentos,
        Configuración, Chat de prueba, Usuarios) y cierre de sesión.
  \item \textbf{Área principal:} contenido de la sección activa.
\end{itemize}

\begin{figure}[H]
  \centering
  \includegraphics[width=0.92\linewidth]{06-shell-navegacion.png}
  \caption{Estructura general del panel con menú lateral}
  \label{fig:documentos-subida}
\end{figure}

Para salir, use \botonnavy{Cerrar sesión} al pie del menú.

% ============================================================
\chapter{Inicio}

\section{Descripción}
La sección \menupath{Inicio} muestra el estado general del asistente:

\begin{itemize}
  \item Estado del servicio.
  \item Proveedor de IA activo (NVIDIA o Gemini).
  \item Modelo LLM en uso.
  \item Cantidad de documentos en la biblioteca.
  \item Aviso si hace falta \textbf{reindexar}.
\end{itemize}

\begin{figure}[H]
  \centering
  \includegraphics[width=0.92\linewidth]{07-dashboard.png}
  \caption{Tablero de inicio con métricas del servicio}
  \label{fig:documentos-subida}
\end{figure}

\begin{notabox}
  Si aparece un aviso de reindexación pendiente, vaya a \menupath{Documentos}
  y ejecute \botonnavy{Reindexar} antes de confiar en el chat público o de prueba.
\end{notabox}

% ============================================================
\chapter{Documentos}

\section{Descripción}
En \menupath{Documentos} gestiona el corpus que alimenta al asistente:
subida, eliminación y \\
actualización del índice de búsqueda (reindexación).

\begin{tipobox}
  Flujo recomendado: \\
  \textbf{1)} Subir documento \\
  \textbf{2)} Reindexar \\
  \textbf{3)} Probar en Chat de prueba.
\end{tipobox}

\section{Formatos admitidos}
\begin{itemize}
  \item PDF (.pdf)
  \item Texto plano (.txt)
  \item Markdown (.md)
  \item Tamaño máximo configurado en el servidor (por defecto 25~MB).
\end{itemize}

\section{Subir un documento}

\begin{pasobox}{1}
  Entre a \menupath{Documentos}.
\end{pasobox}
\begin{pasobox}{2}
  Use el área de carga (ícono de subida) o el selector de archivo para elegir
  el documento.
\end{pasobox}
\begin{pasobox}{3}
  Confirme con \boton{Subir documento} (o el botón equivalente con ícono upload).
\end{pasobox}

\begin{figure}[H]
  \centering
  \includegraphics[width=0.92\linewidth]{08-documentos-subida.png}
  \caption{Sección de carga de archivos}
  \label{fig:documentos-subida}
\end{figure}

\begin{advertenciabox}
  Tras subir o eliminar un archivo, el índice queda desactualizado hasta que
  pulse \textbf{Reindexar}. El tablero mostrará el aviso correspondiente.
\end{advertenciabox}

\section{Biblioteca de archivos}
La tabla lista nombre, tamaño y fecha de actualización. Puede abrir el archivo
o eliminarlo con el ícono de basurero (se pedirá confirmación).

\begin{figure}[h]
  \centering
  \includegraphics[width=0.92\linewidth]{09-documentos-biblioteca.png}
  \caption{Biblioteca de documentos con acciones}
  \label{fig:documentos-subida}
\end{figure}

\section{Reindexar}

\begin{pasobox}{1}
  Pulse \botonnavy{Reindexar} (ícono de actualización).
\end{pasobox}
\begin{pasobox}{2}
  Confirme en el diálogo. El proceso puede tardar según el volumen de documentos
  y el proveedor de embeddings.
\end{pasobox}
\begin{pasobox}{3}
  Al finalizar, verá un mensaje de éxito con el resumen de documentos/fragmentos indexados.
\end{pasobox}

\begin{figure}[H]
  \centering
  \includegraphics[width=0.92\linewidth]{10-documentos-reindex.png}
  \caption{Confirmación de reindexación}
  \label{fig:documentos-subida}
\end{figure}

\begin{notabox}
  Solo los administradores autenticados pueden reindexar. No existe un endpoint
  público de indexación.
\end{notabox}

% ============================================================
\chapter{Configuración}

\section{Descripción}
\menupath{Configuración} agrupa tres bloques:
\begin{enumerate}
  \item Claves API.
  \item Modelo de IA.
  \item Personalización visual.
\end{enumerate}

\section{Claves API}
Permite registrar las claves de \textbf{Google Gemini} y \textbf{NVIDIA}.

\begin{itemize}
  \item Se almacenan \textbf{cifradas} en la base de datos.
  \item Si deja el campo vacío al guardar, \textbf{no se modifica} la clave actual.
  \item Prioridad: clave en base de datos $\rightarrow$ variable de entorno (.env).
\end{itemize}

\begin{figure}[H]
  \centering
  \includegraphics[width=0.92\linewidth]{12-config-claves.png}
  \caption{Sección de claves API con información de estado}
  \label{fig:documentos-subida}
\end{figure}

\begin{advertenciabox}
  NVIDIA es necesaria para los \textbf{embeddings} (incluso si el LLM es Gemini).
  Sin ella no se puede reindexar ni recuperar fragmentos correctamente.
\end{advertenciabox}

\section{Modelo de IA}
Campos principales:

\begin{tabularx}{\linewidth}{L{4.2cm} X}
  \toprule
  \rowcolor{navy}
  \textcolor{white}{\textbf{Campo}} &
  \textcolor{white}{\textbf{Descripción}} \\
  \midrule
  Proveedor LLM & NVIDIA o Google Gemini. \\
  \rowcolor{navysoft}
  Modelo LLM & Identificador del modelo (ej. meta/llama-3.1-8b-instruct). \\
  Modelo de embeddings & Modelo NVIDIA para vectorizar documentos. \\
  \rowcolor{navysoft}
  Fragmentos por consulta & Entre 1 y 20 (cuántos trozos de documento recibe el modelo). \\
  \bottomrule
\end{tabularx}

\begin{figure}[H]
  \centering
  \includegraphics[width=0.92\linewidth]{13-config-modelo.png}
  \caption{Formulario de proveedor, modelo y fragmentos}
  \label{fig:documentos-subida}
\end{figure}

\begin{advertenciabox}
  Si cambia el \textbf{modelo de embeddings}, el índice queda desactualizado:
  debe reindexar obligatoriamente.
\end{advertenciabox}

\section{Personalización}

\begin{itemize}
  \item \campo{URL del logo} (preferible PNG/SVG horizontal con fondo transparente).
  \item \campo{Nombre} y \campo{Subtítulo}.
  \item \campo{Color primario}.
  \item \campo{Color de acento}.
\end{itemize}

\begin{figure}[H]
  \centering
  \includegraphics[width=0.92\linewidth]{14-config-branding.png}
  \caption{Personalización: logo, nombre y colores}
  \label{fig:documentos-subida}
\end{figure}

Pulse \boton{Guardar configuración} al finalizar. Los cambios de branding se
aplican de inmediato en el panel.

% ============================================================
\chapter{Chat de prueba}

\section{Descripción}
\menupath{Chat de prueba} envía preguntas al motor RAG con la configuración
actual, sin pasar por el widget público. Es el lugar ideal para validar
documentos recién indexados.

\begin{pasobox}{1}
  Escriba una pregunta representativa de la normativa o del documento a validar.
\end{pasobox}
\begin{pasobox}{2}
  Pulse el botón de envío (ícono de papel).
\end{pasobox}
\begin{pasobox}{3}
  Revise la respuesta y, si aparecen, las fuentes asociadas.
\end{pasobox}

\begin{figure}[H]
  \centering
  \includegraphics[width=0.92\linewidth]{15-chat-prueba.png}
  \caption{Chat de prueba con pregunta y respuesta}
  \label{fig:documentos-subida}
\end{figure}

\begin{tipobox}
  Ejemplo de pregunta de control:\\
  \textit{«¿Qué pasa si pierdo más de 3 veces un curso?»}
\end{tipobox}


% ============================================================
\chapter{Usuarios}

\section{Descripción}
\menupath{Usuarios} permite crear administradores del panel y activarlos o
desactivarlos.

\section{Crear un administrador}

\begin{pasobox}{1}
  Complete \campo{Nombre completo}, \campo{Correo} y \campo{Contraseña} (mín. 8).
\end{pasobox}
\begin{pasobox}{2}
  Pulse \boton{Crear usuario}.
\end{pasobox}

\begin{figure}[H]
  \centering
  \includegraphics[width=0.92\linewidth]{16-usuarios-crear.png}
  \caption{Formulario de nuevo administrador}
  \label{fig:documentos-subida}
\end{figure}

\section{Listado y estados}
Cada fila muestra nombre, correo y estado ( Activo / Inactivo).
Puede desactivar o activar según corresponda.

\begin{figure}[H]
  \centering
  \includegraphics[width=0.92\linewidth]{17-usuarios-listado.png}
  \caption{Listado de administradores}
  \label{fig:documentos-subida}
\end{figure}

\begin{advertenciabox}
  No es posible desactivar al \textbf{único} administrador activo desde la
  interfaz: el sistema protege el acceso al panel.
\end{advertenciabox}

% ============================================================
\chapter{Preguntas frecuentes y solución de \\ problemas}

\section{Preguntas frecuentes}

\subsection*{¿El chat inventa respuestas?}
No debería: el prompt obliga a basarse solo en documentos recuperados. Si no hay
evidencia, responde que no tiene la información.

\subsection*{Subí un PDF pero el chat no lo menciona}
Verifique que ejecutó \textbf{Reindexar} después de subirlo y pruebe en
\menupath{Chat de prueba}.

\subsection*{Cambié a Gemini y falló}
Confirme que existe clave Gemini (BD o .env) y que NVIDIA sigue
disponible para embeddings.

\subsection*{Olvidé la contraseña de admin}
Otro administrador puede crear un usuario nuevo o actualizar la contraseña
desde el backend/seed. No hay recuperación automática por correo en el MVP.

\section{Errores comunes}
\begin{tabularx}{\linewidth}{L{4.5cm} X}
  \toprule
  \rowcolor{navy}
  \textcolor{white}{\textbf{Situación}} &
  \textcolor{white}{\textbf{Qué hacer}} \\
  \midrule
  No inicia sesión & Verifique correo, contraseña y que el usuario esté activo. \\
  \rowcolor{navysoft}
  Fallo al subir archivo & Revise formato (PDF/TXT/MD) y tamaño máximo. \\
  Reindex falla & Revise clave NVIDIA y que haya al menos un documento válido. \\
  \rowcolor{navysoft}
  Widget no aparece & En WordPress: widget habilitado y plugin activo. \\
  \bottomrule
\end{tabularx}

% ============================================================
\chapter{Glosario}

\begin{description}[style=nextline, leftmargin=2.2cm, labelindent=0.3cm]
  \item[\textcolor{navy}{RAG}]
    \textit{Retrieval-Augmented Generation}: recupera fragmentos de documentos
    y genera la respuesta con un modelo de lenguaje.
  \item[\textcolor{navy}{Embeddings}]
    Representaciones numéricas del texto usadas para búsqueda semántica.
  \item[\textcolor{navy}{Reindexar}]
    Volver a procesar los documentos y actualizar el índice vectorial.
  \item[\textcolor{navy}{LLM}]
    \textit{Large Language Model}: modelo generativo (NVIDIA NIM o Gemini).
  \item[\textcolor{navy}{Widget}]
    Componente de chat flotante en el sitio WordPress.
  \item[\textcolor{navy}{JWT}]
    Token de autenticación del panel de administración.
\end{description}

\end{document}
